\scanpage{251}
\mspage{201}

semblent n'y jouer pratiquement aucun rôle, ou un rôle très effacé.

Le carré de la p.~197 montre surtout que la catégorie $\unsure{\mathcal C|\mathscr C}$ est filtrante -- $(C,i),(C',i')$ s'envoient l'un et l'autre dans $(C'',i'')$, avec $C''=C\circ C'$, \deleted{\unsure{Alors}} $i''=(C*i')\circ i=(i*C')\circ i'$. Il y a $=$

une loi de composition dans $\mathcal C$, nullement commutative ni unitaire c'est sûr, mais sûrement associative. et \unsure{(}\deleted{\unsure{soit|si}} cette composition est notée $*$) avec des morphismes fonctoriels

\diagram{\begin{tikzcd}[ampersand replacement=\&]
{\unsure{\xi|\zeta}}\arrow[r,"{\alpha_{\xi,\xi'}}"]\&\xi*\xi'\&\xi'\arrow[l,"{\beta_{\xi,\xi'}}"']
\end{tikzcd}}

qui incarnent le caractère filtrant croissant de $\mathcal C$. Ces flèches semblent \unsure{remplacer} en quelque sorte l'unité bilatère absente, et on s'attend à ce qu'elles aient les mêmes excellentes propriétés que si on avait \unsure{une} telle unité, concernant les composés doubles $((\xi\unsure{\circ|*}\xi')\unsure{\times|*}\xi''\simeq\xi\circ(\xi'\unsure{\circ|*}\xi''\unsure{)}$ \unsure{:}
% archive sheet number: 251
